% Problem set with solutions. The `answers` option prints the solutions.
% Remove `answers` to hide them and print a blank assignment.
\documentclass[11pt,addpoints,answers]{exam}

\usepackage[T1]{fontenc}
\usepackage{amsmath,amssymb}
\usepackage[margin=1in]{geometry}

% Header and footer. The first page has its own header.
\pagestyle{headandfoot}
\firstpageheader{Course 101}{Problem Set 3}{Due: Friday, 11:59 pm}
\runningheader{Course 101}{Problem Set 3}{Due: Friday, 11:59 pm}
\firstpagefooter{}{Page \thepage}{}
\runningfooter{}{Page \thepage}{}

\begin{document}

\begin{center}
  {\Large\bfseries {{title}}}\\[0.6em]
  Student: {{author}}
\end{center}

\vspace{0.5em}

\begin{questions}

  % Parts take their own points; the question total is the sum of its parts.
  \question[10] Let $f(x) = x^3 - 6x^2 + 9x$.
  \begin{parts}
    \part[4] Find the critical points of $f$.
    \begin{solution}
      $f'(x) = 3x^2 - 12x + 9 = 3(x-1)(x-3)$, so the critical points are $x = 1$ and $x = 3$.
    \end{solution}

    \part[3] Classify each critical point as a local maximum or a local minimum.
    \begin{solution}
      $f''(x) = 6x - 12$. Since $f''(1) = -6 < 0$, $x = 1$ is a local maximum.
      Since $f''(3) = 6 > 0$, $x = 3$ is a local minimum.
    \end{solution}

    \part[3] Explain why $f$ has no absolute maximum on the whole real line.
    \begin{solution}
      As $x \to \infty$, $f(x) \to \infty$, so $f$ is unbounded above and no
      value of $f$ is largest.
    \end{solution}
  \end{parts}

  % Multiple choice: \CorrectChoice marks the answer shown in the solutions.
  \question[5] Which statement about the sample variance $s^2$ is correct?
  \begin{choices}
    \choice It divides the sum of squared deviations by $n$.
    \CorrectChoice It divides the sum of squared deviations by $n-1$, which makes it an unbiased estimator of $\sigma^2$.
    \choice It is unaffected by a single extreme value.
    \choice It is always negative.
  \end{choices}
  \begin{solution}
    The denominator $n-1$ corrects for the fact that the deviations are
    measured from the sample mean, which uses up one degree of freedom.
  \end{solution}

  \question[12] Calls to a support line arrive as a Poisson process at a rate of
  3 per minute. Let $X$ be the number of calls in one minute.
  \begin{parts}
    \part[4] Compute $P(X = 0)$.
    \begin{solution}
      $P(X = 0) = e^{-3} \approx 0.0498$.
    \end{solution}

    \part[4] Compute $P(X \geq 2)$.
    \begin{solution}
      $P(X \geq 2) = 1 - P(X = 0) - P(X = 1) = 1 - 4e^{-3} \approx 0.801$.
    \end{solution}

    \part[4] Find the mean and the variance of the number of calls in a two-minute interval.
    \begin{solution}
      The count in two minutes is Poisson with mean $6$, so the mean and the
      variance are both $6$.
    \end{solution}
  \end{parts}

  \question[10] Prove that the sum of two even integers is even. Write your
  proof in full sentences.
  \begin{solution}
    Let $a = 2m$ and $b = 2n$ for integers $m$ and $n$. Then
    $a + b = 2(m+n)$, which is twice the integer $m+n$, so $a+b$ is even.
  \end{solution}

  % Bonus points are counted separately from the regular total.
  \bonusquestion[5] For a Poisson random variable $X$ with mean $\lambda$, show that $E[X] = \lambda$.
  \begin{solution}
    \[
      E[X] = \sum_{k=0}^{\infty} k\,\frac{e^{-\lambda}\lambda^k}{k!}
      = \lambda e^{-\lambda} \sum_{k=1}^{\infty} \frac{\lambda^{k-1}}{(k-1)!}
      = \lambda e^{-\lambda} e^{\lambda} = \lambda.
    \]
  \end{solution}

\end{questions}

\vspace{1em}

\gradetable[h][questions]

\end{document}
